\documentclass[11pt,a4paper]{article}

% ── Packages ──────────────────────────────────────────────────────────────────
\usepackage[margin=2.5cm]{geometry}
\usepackage{amsmath,amssymb,amsthm}
\usepackage{booktabs}
\usepackage{graphicx}
\usepackage{hyperref}
\usepackage{enumitem}
\usepackage{algorithm}
\usepackage{algpseudocode}
\usepackage{xcolor}
\usepackage{listings}
\usepackage{multirow}
\usepackage{longtable}
\usepackage{caption}
\usepackage{subcaption}
\usepackage{fancyhdr}

\hypersetup{
  colorlinks=true,
  linkcolor=blue!60!black,
  citecolor=green!50!black,
  urlcolor=blue!70!black,
}

\lstset{
  language=Python,
  basicstyle=\ttfamily\small,
  keywordstyle=\color{blue!70!black},
  commentstyle=\color{green!50!black}\itshape,
  stringstyle=\color{red!60!black},
  frame=single,
  breaklines=true,
  numbers=left,
  numberstyle=\tiny\color{gray},
}

\pagestyle{fancy}
\fancyhead[L]{\small Polar Codes for the Two-User MAC}
\fancyhead[R]{\small\thepage}
\fancyfoot{}

% ── Math shortcuts ────────────────────────────────────────────────────────────
\newcommand{\GF}{\mathrm{GF}}
\newcommand{\ZZ}{\mathcal{Z}}
\DeclareMathOperator{\logsumexp}{logsumexp}
\DeclareMathOperator{\circconv}{\circledast}

% ══════════════════════════════════════════════════════════════════════════════
\begin{document}

\begin{titlepage}
\centering
\vspace*{3cm}
{\Huge\bfseries Polar Codes for the\\[0.3em] Two-User Multiple Access Channel}\\[1.5cm]
{\Large Design, Decoding, and Neural Decoder Approaches}\\[2cm]
{\large Project Documentation}\\[1cm]
{\large March 2026}\\[3cm]
\begin{abstract}
\noindent
This document provides a comprehensive description of a software library for
polar codes on the two-user binary-input multiple access channel (MAC).
The project supports three channel models---Binary Erasure MAC (BEMAC),
Additive Binary Noise MAC (ABNMAC), and Gaussian MAC---with both
analytical and Monte Carlo code design.
Five decoder implementations are provided, ranging from an $O(N^2)$
reference decoder to an $O(N \log N)$ computational graph decoder
supporting all monotone chain decoding paths, a Numba-JIT parallel
decoder, and an SC List (SCL) decoder.
The project also includes a neural decoder based on the computational
graph framework, using learned MLP operations for tree operations
including a fully neural CalcParent.
Extensive simulation results (BLER vs.\ block length, SNR sweeps, rate
analysis) are presented for all channel models and code classes.
\end{abstract}
\end{titlepage}

\tableofcontents
\newpage

% ══════════════════════════════════════════════════════════════════════════════
\section{Executive Summary}
\label{sec:executive}

This project implements polar codes for the two-user binary-input multiple
access channel (MAC), following the framework of \"{O}nay~\cite{onay2013}
and Ren, Nasser, and Bhatt~\cite{ren2025}.

\subsection{Key Contributions}

\begin{enumerate}[label=(\arabic*)]
  \item \textbf{Three channel models:} BEMAC ($Z = X + Y$), ABNMAC
        ($Z = (X \oplus E_x, Y \oplus E_y)$), and Gaussian MAC
        ($Z = (1{-}2X) + (1{-}2Y) + W$).

  \item \textbf{Analytical and Monte Carlo code design:}
        Bhattacharyya-parameter recursion (Arikan 2009) for BEMAC and ABNMAC;
        genie-aided MC design for arbitrary channels and paths.

  \item \textbf{Five decoder backends:}
        \begin{itemize}
          \item Reference $O(N^2)$ SC decoder (\texttt{decoder.py})
          \item Efficient $O(N \log N)$ LLR-based decoder for extreme paths (\texttt{efficient\_decoder.py})
          \item Interleaved $O(N \log N)$ computational graph decoder for all
                monotone chain paths (\texttt{decoder\_interleaved.py})
          \item Numba-JIT parallel decoder with batch \texttt{prange} parallelism
                (\texttt{decoder\_parallel.py})
          \item SCL decoder with list sizes up to $L = 32$
                (\texttt{decoder\_scl.py})
        \end{itemize}

  \item \textbf{Neural SC decoder (NCG):}
        MLP-based tree operations with a fully neural CalcParent,
        matching or beating analytical SC for BEMAC Class B at $N = 8$--$128$.

  \item \textbf{Extensive simulation results:}
        BEMAC BLER vs.\ $N$ for Classes A, B, C; ABNMAC across block lengths
        $N = 8$--$1024$; Gaussian MAC campaign with 824 data points
        (3 classes $\times$ multiple SNRs $\times$ $N = 8$--$2048$, SC and SCL).
\end{enumerate}

\subsection{Key Findings}

\begin{itemize}
  \item \textbf{BEMAC Class C} (path $0^N 1^N$) achieves BLER $\approx 2.6 \times 10^{-4}$ at $N = 1024$ with SC decoding at rate ratio $\rho = 0.6$.
  \item \textbf{BEMAC Class A} (symmetric rate direction) suffers from error floor: BLER $\to 1.0$ as $N$ increases, confirming that the symmetric path is capacity-suboptimal for the asymmetric BEMAC capacity region.
  \item \textbf{Gaussian MAC:} SC performance is strongly SNR-dependent. At $\text{SNR} = 6$ dB and rate ratio $\rho = 0.5$, Classes A and B achieve $\text{BLER} \to 0$ as $N$ grows, while $\rho = 0.9$ causes error floors.
  \item \textbf{NCG neural decoder:} At $N = 128$ on BEMAC Class B, the neural decoder achieves BLER $= 0.0018$ vs.\ SC BLER $= 0.0034$---a $47\%$ improvement.
\end{itemize}


% ══════════════════════════════════════════════════════════════════════════════
\section{Background and Theory}
\label{sec:background}

\subsection{Polar Codes (Arikan 2009)}

Polar codes are a family of capacity-achieving codes based on the phenomenon
of \emph{channel polarization}~\cite{arikan2009}. Given a binary-input
memoryless channel $W$, one constructs $N = 2^n$ synthetic sub-channels
$W_N^{(i)}$ for $i = 1, \ldots, N$. As $N \to \infty$, each sub-channel
polarizes to either a noiseless channel or a completely noisy channel.

The generator matrix is $G_N = B_N \cdot F^{\otimes n}$ where $F = \begin{pmatrix} 1 & 0 \\ 1 & 1 \end{pmatrix}$ is the Arikan kernel and $B_N$ is the
bit-reversal permutation matrix. Encoding is $x^N = u^N \cdot G_N$ over
$\GF(2)$.

The Bhattacharyya parameter $Z(W) \in [0,1]$ measures channel reliability.
The Arikan recursion for polarization is:
\begin{align}
  Z(W^-) &= 2 Z(W) - Z(W)^2 \quad \text{(bad split, higher $Z$)} \\
  Z(W^+) &= Z(W)^2 \quad \text{(good split, lower $Z$)}
\end{align}
Information bits are placed on sub-channels with $Z(W_N^{(i)}) \approx 0$
(reliable); the remaining positions are set to known frozen values.

\subsection{Multiple Access Channel (MAC) Model}

The two-user binary-input MAC has inputs $X \in \{0,1\}$ (User 1) and
$Y \in \{0,1\}$ (User 2), with channel transition probability
$W(z | x, y)$. The capacity region under uniform input distributions is
characterized by:
\begin{align}
  R_x &\leq I(Z; X) \\
  R_y &\leq I(Z; Y | X) \\
  R_x + R_y &\leq I(Z; X, Y)
\end{align}

The polar coding approach for the MAC treats the two users' messages as
two separate polar codes that share the same channel output $z^N$.
Successive cancellation (SC) decoding alternates between decoding User 1
(U) bits and User 2 (V) bits according to a \emph{decoding path} $b^{2N}$.

\subsection{Channel Models}

\subsubsection{Binary Erasure MAC (BEMAC)}

\begin{equation}
  Z = X + Y \in \{0, 1, 2\}
\end{equation}
This is a deterministic channel. The output $Z = 1$ is ambiguous (either
$(X,Y) = (0,1)$ or $(1,0)$), while $Z = 0$ and $Z = 2$ uniquely determine
both inputs. The capacity region (uniform inputs) is:
\[
  R_x \leq 0.5, \quad R_y \leq 0.5, \quad R_x + R_y \leq 1.0
\]

The Bhattacharyya starting values for the extreme path $0^N 1^N$ are:
\begin{itemize}
  \item $Z_0^{(U)} = 0.5$ (U marginal effective channel is BEC(0.5))
  \item $Z_0^{(V)} = 0.0$ (V conditional channel is noiseless once $X$ is known)
\end{itemize}

\subsubsection{Additive Binary Noise MAC (ABNMAC)}

\begin{equation}
  Z = (X \oplus E_x, \; Y \oplus E_y), \quad (E_x, E_y) \sim p_{\text{noise}}
\end{equation}
The noise distribution is a $2 \times 2$ joint probability matrix. The
default parameters (from \"{O}nay 2013, Table I) are:
\[
  p_{\text{noise}} = \begin{pmatrix} 0.1286 & 0.0175 \\ 0.0175 & 0.8364 \end{pmatrix}
\]
yielding capacities $I(Z;X) \approx 0.400$, $I(Z;Y|X) \approx 0.800$,
$I(Z;X,Y) \approx 1.200$ bits. The Bhattacharyya starting values are:
\begin{align}
  Z_0^{(U)} &= 2\sqrt{p_{E_x}(0) \cdot p_{E_x}(1)} \approx 0.706 \\
  Z_0^{(V)} &= 2\bigl[\sqrt{p_{00} p_{01}} + \sqrt{p_{10} p_{11}}\bigr] \approx 0.337
\end{align}

\subsubsection{Gaussian MAC}

\begin{equation}
  Z = (1 - 2X) + (1 - 2Y) + W, \quad W \sim \mathcal{N}(0, \sigma^2)
\end{equation}
where $X, Y \in \{0,1\}$ are BPSK-modulated: $s_x = 1 - 2X \in \{-1, +1\}$.
The per-user SNR is $\text{SNR} = 1/\sigma^2$. The noiseless constellation
points are $\{-2, 0, 0, +2\}$ for $(X,Y) \in \{(1,1), (0,1), (1,0), (0,0)\}$.

The V-conditional Bhattacharyya parameter has a closed form:
$Z_0^{(V)} = e^{-1/(2\sigma^2)}$, while the U-marginal parameter
$Z_0^{(U)}$ requires numerical integration over the mixture Gaussian
marginal channel.

\subsection{Code Classes and Decoding Paths}

The decoding path $b^{2N}$ is a binary vector of length $2N$ that specifies
the order in which U and V bits are decoded. A $0$ in position $k$ means
``decode one U bit''; a $1$ means ``decode one V bit.'' The project uses
the \emph{monotone chain} path family:
\begin{equation}
  b^{2N} = 0^{\pi} \, 1^N \, 0^{N - \pi}, \quad 0 \leq \pi \leq N
\end{equation}
Key code classes are:
\begin{itemize}
  \item \textbf{Class C} ($\pi = N$): $b = 0^N 1^N$ --- decode all U bits first, then all V bits. Corresponds to $R_U$ direction along the U-axis of the capacity region.
  \item \textbf{Class A} ($\pi = 0$): $b = 1^N 0^N$ --- decode all V bits first, then all U bits.
  \item \textbf{Class B} ($\pi = N/2$): Intermediate path, interleaving U and V decoding.
\end{itemize}

For extreme paths (Classes A and C), the decoder decomposes into two
sequential single-user SC decodes with $O(N \log N)$ complexity each.
For intermediate paths, the computational graph approach~\cite{ren2025}
is required.


% ══════════════════════════════════════════════════════════════════════════════
\section{Code Design}
\label{sec:design}

\subsection{Analytical Bhattacharyya Design (\texttt{design.py})}

The analytical design applies $n$ stages of the Arikan recursion starting
from a channel-specific initial Bhattacharyya parameter $Z_0$.

\begin{algorithm}[h]
\caption{Bhattacharyya Recursion}
\begin{algorithmic}[1]
\Require Starting parameter $Z_0 \in [0,1]$, stages $n$
\State $z \gets [Z_0]$
\For{stage $= 1, \ldots, n$}
  \State $z_{\text{bad}} \gets \min(1, 2z - z^2)$ \Comment{$Z^-$ split}
  \State $z_{\text{good}} \gets z^2$ \Comment{$Z^+$ split}
  \State Interleave: $z_{\text{new}}[2k] = z_{\text{bad}}[k]$, $z_{\text{new}}[2k{+}1] = z_{\text{good}}[k]$
\EndFor
\State \Return $z$ of length $N = 2^n$
\end{algorithmic}
\end{algorithm}

For each channel model, the starting parameters are computed analytically:
\begin{itemize}
  \item \textbf{BEMAC:} $Z_0^{(U)} = 0.5$, $Z_0^{(V)} = 0.0$
  \item \textbf{ABNMAC:} $Z_0^{(U)} = 2\sqrt{p_{E_x}(0) \cdot p_{E_x}(1)}$,
        $Z_0^{(V)} = 2[\sqrt{p_{00} p_{01}} + \sqrt{p_{10} p_{11}}]$
\end{itemize}

The best $k_u$ (resp.\ $k_v$) sub-channels with smallest Bhattacharyya
parameters form the information sets $A_U$ (resp.\ $A_V$). All positions
use 1-based indexing throughout the codebase.

The function \texttt{design\_bemac(n, ku, kv)} returns:
\begin{itemize}
  \item $A_U$, $A_V$: sorted lists of 1-indexed information positions
  \item \texttt{frozen\_u}, \texttt{frozen\_v}: dictionaries mapping frozen positions to their values (typically 0)
  \item $z_u$, $z_v$: Bhattacharyya parameter arrays
\end{itemize}

\subsection{Monte Carlo Genie-Aided Design (\texttt{design\_mc.py})}

For channels where analytical Bhattacharyya computation is unavailable
(e.g., Gaussian MAC) or for intermediate paths, the project provides a
Monte Carlo design procedure:

\begin{enumerate}
  \item Generate random message pairs $(u, v)$
  \item Encode both users: $x = u \cdot G_N$, $y = v \cdot G_N$
  \item Transmit through the channel to get $z$
  \item Run a \emph{genie-aided} SC decoder that always feeds the true bit value to subsequent steps, but records whether each bit would have been decoded incorrectly
  \item Accumulate per-bit error counts over many trials
  \item Rank sub-channels by error rate (ascending = most reliable first)
\end{enumerate}

Tie-breaking uses bit-reversal weight: when multiple channels have identical
MC error rates (common when $P_e = 0$ for many channels), the natural polar
reliability ordering breaks ties via the bit-reversed index.

The design files are stored as NumPy \texttt{.npz} archives containing
\texttt{sorted\_u}, \texttt{sorted\_v}, \texttt{error\_rates\_u},
\texttt{error\_rates\_v}, and optionally \texttt{path\_i}.

\subsection{Gaussian MAC Design (\texttt{design\_mc.py})}

The Gaussian MAC presents a unique design challenge: the continuous output
alphabet makes analytical Bhattacharyya computation more complex. While the
V-conditional parameter has a closed-form $Z_0^{(V)} = e^{-1/(2\sigma^2)}$,
the U-marginal parameter requires numerical integration over the mixture
Gaussian channel.

For non-extreme paths (Classes A and B), the analytical Bhattacharyya approach
does not directly apply because the effective channels depend on the decoding
order. The MC design procedure handles this naturally by simulating the
actual genie-aided SC decoder along the specified path.

The Gaussian MAC design workflow is:
\begin{enumerate}
  \item Choose class (A, B, or C), block length $N = 2^n$, and SNR in dB
  \item Create a \texttt{GaussianMAC} channel: $\sigma^2 = 10^{-\text{SNR}_{\text{dB}}/10}$
  \item Run \texttt{mc\_design(n, channel, mc\_trials, path\_i=pi)} to
        estimate per-bit error rates for both users
  \item Select the $k_u$ and $k_v$ most reliable sub-channels
  \item Save the design to \texttt{designs/gmac\_\{class\}\_n\{n\}\_snr\{snr\}dB.npz}
\end{enumerate}

The MC design uses a time budget to automatically stop after a specified
number of seconds, making it practical for large block lengths where each
trial is expensive ($O(N^2)$ for the genie-aided decoder).

\subsection{Rate Selection}

Given a target sum rate $R_{\text{sum}}$ and a rate ratio $\rho = R_u / R_{\text{sum}}$,
the number of information bits is:
\begin{align}
  k_u &= \lfloor \rho \cdot R_{\text{sum}} \cdot N \rceil \\
  k_v &= \lfloor (1 - \rho) \cdot R_{\text{sum}} \cdot N \rceil
\end{align}
where $\lfloor \cdot \rceil$ denotes rounding to the nearest integer.
The actual rates $(k_u/N, k_v/N)$ are thus quantized, and the gap between
target and actual rates can be significant at small $N$.

Additionally, $k_u$ and $k_v$ must not exceed the number of reliable
sub-channels: $k_u \leq |\{i : Z_u(i) < \text{threshold}\}|$ and similarly
for $k_v$. The project uses MC error rates rather than Bhattacharyya
parameters for this assessment when MC designs are available.

\subsection{Design File Inventory}

The \texttt{designs/} directory contains pre-computed designs:
\begin{itemize}
  \item \textbf{BEMAC:} Classes A, B, C for $n = 3, \ldots, 10$ ($N = 8, \ldots, 1024$). Additional $n = 12$ for Class B.
  \item \textbf{ABNMAC:} Classes A, B, C for $n = 3, \ldots, 10$.
  \item \textbf{Gaussian MAC:} Classes A, B, C for $n = 3, \ldots, 11$ at SNRs $\{0, 1, 2, 3, 4, 5, 6, 8, 10\}$ dB. Over 250 design files total.
\end{itemize}


% ══════════════════════════════════════════════════════════════════════════════
\section{Encoder (\texttt{encoder.py})}
\label{sec:encoder}

The polar encoder implements $x^N = u^N \cdot G_N = u^N \cdot B_N \cdot F^{\otimes n}$ over $\GF(2)$.

\subsection{Implementation}

The encoding proceeds in two stages:
\begin{enumerate}
  \item \textbf{Bit-reversal permutation:} Apply $B_N$ by reordering the input vector according to the bit-reversal of each index. For index $i$ with $n$-bit representation, compute $\text{br}[i] = \text{reverse\_bits}(i, n)$, then set $x[\text{br}[i]] = u[i]$.

  \item \textbf{XOR butterfly:} Apply $F^{\otimes n}$ via $n$ stages of butterfly XOR operations. At each stage with step size $s = 1, 2, \ldots, N/2$, the array is viewed as blocks of size $2s$, and the left half of each block is XORed with the right half:
  \begin{equation}
    x_r[:, :s] \mathrel{\oplus}= x_r[:, s:]
  \end{equation}
\end{enumerate}

The complexity is $O(N \log N)$ per codeword, or $O(B \cdot N \log N)$ for batch encoding of $B$ codewords simultaneously.

\subsection{Batch and TensorFlow Variants}

Three encoder variants are provided:
\begin{itemize}
  \item \texttt{polar\_encode(u)}: Single-vector encoding, returns a list.
  \item \texttt{polar\_encode\_batch(U)}: Batch NumPy encoding of shape $(B, N)$.
  \item \texttt{polar\_encode\_batch\_tf(U, N)}: TensorFlow GPU-accelerated batch encoding using \texttt{tf.bitwise.bitwise\_xor}.
\end{itemize}

The helper \texttt{build\_message(N, info\_bits, info\_positions)} places
information bits at the specified 1-indexed positions, with all other
positions set to their frozen values (default: 0).


% ══════════════════════════════════════════════════════════════════════════════
\section{Decoders}
\label{sec:decoders}

\subsection{Overview}

The project provides five decoder implementations with increasing
sophistication:

\begin{table}[h]
\centering
\caption{Decoder summary}
\begin{tabular}{llll}
\toprule
\textbf{Decoder} & \textbf{Complexity} & \textbf{Paths} & \textbf{File} \\
\midrule
Reference SC  & $O(N^2)$ & All & \texttt{decoder.py} (legacy) \\
Efficient SC  & $O(N \log N)$ & Extreme only & \texttt{efficient\_decoder.py} \\
Interleaved SC & $O(N \log N)$ & All monotone chain & \texttt{decoder\_interleaved.py} \\
Parallel SC   & $O(N \log N)$ & Class B & \texttt{decoder\_parallel.py} \\
SC List (SCL) & $O(L \cdot N \log N)$ & All & \texttt{decoder\_scl.py} \\
\bottomrule
\end{tabular}
\end{table}

The unified public API in \texttt{decoder.py} auto-dispatches:
\begin{itemize}
  \item Extreme paths ($\pi = 0$ or $\pi = N$) $\to$ LLR-based SC
  \item Class B ($\pi = N/2$) $\to$ Numba parallel decoder (if available)
  \item General paths $\to$ Tensor-based computational graph SC
\end{itemize}

\subsection{LLR-Based SC Decoder (Extreme Paths)}

For the extreme path $0^N 1^N$ (Class C), the MAC decode decomposes into:
\begin{enumerate}
  \item Compute U-marginal leaf LLRs by marginalizing over $Y$:
  \begin{equation}
    \Lambda_U(t) = \ln \frac{P(z_t | x{=}0)}{P(z_t | x{=}1)} = \ln \frac{\sum_y W(z_t|0,y)}{\sum_y W(z_t|1,y)}
  \end{equation}
  \item Run standard Arikan SC decoding on $\Lambda_U$ to obtain $\hat{u}^N$
  \item Encode $\hat{u}^N$ to get $\hat{x}^N = \hat{u}^N \cdot G_N$
  \item Compute V-conditional leaf LLRs given the decoded $\hat{x}^N$:
  \begin{equation}
    \Lambda_V(t) = \ln \frac{W(z_t | \hat{x}_t, 0)}{W(z_t | \hat{x}_t, 1)}
  \end{equation}
  \item Run standard SC on $\Lambda_V$ to obtain $\hat{v}^N$
\end{enumerate}

The Arikan SC tree uses an $O(N)$ lazy-evaluation structure (\texttt{\_SCNode}
class) with $f$-node (check) and $g$-node (variable) operations:
\begin{align}
  f(\Lambda_a, \Lambda_b) &= \ln(1 + e^{\Lambda_a + \Lambda_b}) - \ln(e^{\Lambda_a} + e^{\Lambda_b}) \label{eq:fnode} \\
  g(\Lambda_a, \Lambda_b, u) &= \Lambda_b + (1 - 2u) \Lambda_a \label{eq:gnode}
\end{align}

\subsection{Computational Graph Decoder (All Paths)}
\label{sec:comp-graph}

For intermediate paths (e.g., Class B with $\pi = N/2$), the decoder uses
the computational graph approach from Ren et al.~\cite{ren2025}. Each edge
in the binary tree carries an array of $2 \times 2$ log-probability tensors
$T[a, b]$ representing the joint distribution over $(X_{\text{local}}, Y_{\text{local}})$.

\subsubsection{Tree Structure and Indexing}

\begin{itemize}
  \item \textbf{Edges:} Numbered $1, \ldots, 2N{-}1$. Edge 1 is the root. Edges $N, \ldots, 2N{-}1$ are leaves.
  \item \textbf{Vertices:} Numbered $1, \ldots, N{-}1$. Vertex $\beta$ has parent edge $\beta$, left child edge $2\beta$, and right child edge $2\beta + 1$.
  \item \textbf{Edge sizes:} Edge $\beta$ at tree level $\ell$ carries $N / 2^\ell$ tensor positions, each a $2 \times 2$ matrix.
  \item \textbf{Leaf mapping:} 1-indexed position $i_t$ maps to leaf edge $i_t + N - 1$ and vertex $(i_t + N - 1) \gg 1$.
\end{itemize}

\subsubsection{Core Operations}

Three operations propagate information through the tree, all operating on
arrays of $(L, 2, 2)$ log-probability tensors:

\paragraph{Circular Convolution ($\circconv$):}
The circular convolution marginalizes over paired positions:
\begin{equation}
  (\circconv(A, B))[a, b] = \ln \sum_{a', b'} \exp\bigl(A[a \oplus a', b \oplus b'] + B[a', b']\bigr)
\end{equation}

\paragraph{Normalized Product ($\odot$):}
Element-wise product (sum in log-domain) followed by normalization:
\begin{equation}
  (A \odot B)[a, b] = A[a,b] + B[a,b] - \ln \sum_{a',b'} \exp(A[a',b'] + B[a',b'])
\end{equation}

\paragraph{calcLeft, calcRight, calcParent:}
\begin{itemize}
  \item \texttt{calcLeft($\beta$):} Given parent edge data and right child, compute left child:
  \begin{equation}
    E_{2\beta} = \circconv\bigl(E_\beta[:l], \; E_\beta[l:] \odot E_{2\beta+1}\bigr)
  \end{equation}
  where $l$ is the size of the child edges.

  \item \texttt{calcRight($\beta$):} Given parent and left child, compute right child:
  \begin{equation}
    E_{2\beta+1} = E_\beta[l:] \odot \circconv\bigl(E_{2\beta}, E_\beta[:l]\bigr)
  \end{equation}

  \item \texttt{calcParent($\beta$):} Given both children, reconstruct parent:
  \begin{equation}
    E_\beta[:l] = \circconv(E_{2\beta}, E_{2\beta+1}), \quad E_\beta[l:] = E_{2\beta+1}
  \end{equation}
\end{itemize}

\subsubsection{Navigation and decodeAt Flow}

The \texttt{step\_to(target)} method navigates from the current vertex to
the target vertex by computing LCA paths and calling \texttt{calcLeft},
\texttt{calcRight}, or \texttt{calcParent} at each step.

The \texttt{decodeAt} procedure for each bit position $i_t$ is:
\begin{enumerate}
  \item Navigate to the target vertex (parent of the leaf edge)
  \item Save the leaf's bottom-up data: $\texttt{temp} = E_{\text{leaf}}[0]$
  \item Compute \texttt{calcLeft} or \texttt{calcRight} to get the top-down message
  \item Combine: $\texttt{combined} = \texttt{temp} \odot E_{\text{leaf}}[0]$
  \item Make hard decision by marginalizing the combined tensor
  \item Update leaf to \emph{partially deterministic}: known dimension $\to$ delta, unknown $\to$ uniform
\end{enumerate}

\subsubsection{DecodeAt Algorithm}

The complete algorithm for decoding a single bit position is:

\begin{algorithm}[h]
\caption{DecodeAt for position $i_t$ at path step $\gamma$ (0=U, 1=V)}
\begin{algorithmic}[1]
\Require $i_t$: 1-indexed position, $\gamma$: user indicator, graph state
\State $\text{leaf\_edge} \gets i_t + N - 1$
\State $\text{target\_vertex} \gets \text{leaf\_edge} \gg 1$
\State \Call{step\_to}{target\_vertex} \Comment{Navigate tree}
\State $\text{temp} \gets E[\text{leaf\_edge}][0].\text{copy}()$ \Comment{Save bottom-up message}
\If{leaf\_edge is even}
  \State \Call{calcLeft}{target\_vertex} \Comment{Computes top-down into leaf}
\Else
  \State \Call{calcRight}{target\_vertex}
\EndIf
\State $\text{combined} \gets \text{norm\_prod}(E[\text{leaf\_edge}][0], \text{temp})$
\If{$i_t$ is frozen}
  \State $\hat{b} \gets \text{frozen\_value}$
\ElsIf{$\gamma = 0$} \Comment{U decision}
  \State $p_0 \gets \text{logsumexp}(\text{combined}[0,0], \text{combined}[0,1])$
  \State $p_1 \gets \text{logsumexp}(\text{combined}[1,0], \text{combined}[1,1])$
  \State $\hat{b} \gets \mathbb{1}[p_1 > p_0]$ \Comment{Marginalize over V}
\Else \Comment{V decision}
  \State $p_0 \gets \text{logsumexp}(\text{combined}[0,0], \text{combined}[1,0])$
  \State $p_1 \gets \text{logsumexp}(\text{combined}[0,1], \text{combined}[1,1])$
  \State $\hat{b} \gets \mathbb{1}[p_1 > p_0]$ \Comment{Marginalize over U}
\EndIf
\State Update leaf to partially deterministic tensor \Comment{Eq.~16 from \cite{ren2025}}
\end{algorithmic}
\end{algorithm}

The partially deterministic leaf update (step 15) is critical for correctness:
after deciding a bit, the leaf tensor must reflect only the known information.
If both U and V values are known at position $i_t$, the leaf becomes a delta:
$T[\hat{u}, \hat{v}] = 1$, all other entries $= 0$ (in log-domain:
$T[\hat{u}, \hat{v}] = 0$, others $= -\infty$). If only one user is known,
the other dimension remains uniform ($\log 0.5$).

\subsubsection{Root Initialization}

A critical implementation detail: the root edge entries must be bit-reversed
from the leaf log-probabilities. Since the encoder applies bit-reversal
before the butterfly, the tree's interleaved pairing structure requires
$\text{root}[t] = \log W[\text{br}[t]]$.

Each root entry is also normalized so that the four entries sum to 1
(i.e., $\sum_{a,b} \exp(T[a,b]) = 1$).

\subsection{Numba Parallel Decoder (\texttt{decoder\_parallel.py})}

The parallel decoder compiles the entire $2N$-step decode loop into Numba
\texttt{@njit} functions with zero Python overhead per decoded bit. Key
features:

\begin{itemize}
  \item \textbf{Flat storage:} All edge tensors are stored in a single contiguous
        \texttt{float64} array with precomputed offsets, avoiding Python object
        overhead.
  \item \textbf{Standalone \texttt{@njit} functions:} Each operation
        (\texttt{\_do\_cl}, \texttt{\_do\_cr}, \texttt{\_do\_cp}, \texttt{\_nav})
        is a standalone compiled function taking explicit array arguments. This
        avoids a Numba closure bug where mutable arrays in closures produce
        silently incorrect results.
  \item \textbf{Batch parallelism:} The \texttt{\_decode\_batch} function uses
        \texttt{prange} to decode multiple codewords in parallel across CPU cores.
  \item \textbf{Performance:} Achieves approximately $20\times$ speedup at
        $N = 1024$ compared to the pure-Python computational graph decoder.
\end{itemize}

\subsection{SC List Decoder (\texttt{decoder\_scl.py})}

The SCL decoder maintains $L$ candidate paths simultaneously, splitting at
each information bit and pruning to keep the $L$ most likely paths. The
implementation provides:

\begin{itemize}
  \item \textbf{Extreme path SCL:} For paths $0^N 1^N$ and $1^N 0^N$, uses
        the Tal--Vardy log-probability propagation~\cite{talvardy2013} with
        vectorized path operations. First decodes one user with SCL, then
        decodes the other user with SCL conditioned on the first.
  \item \textbf{General path SCL:} For intermediate paths, uses the
        computational graph framework with $L$ parallel tree instances and
        path metric tracking.
  \item \textbf{Complexity:} $O(L \cdot N \log N)$ per codeword.
\end{itemize}

The SCL decoder supports list sizes $L = 1, 2, 4, 8, 16, 32$.

\subsection{Batch Vectorization}

All decoders support batch decoding where multiple codewords are processed
in parallel via NumPy vectorization. The batched computational graph
(\texttt{\_CompGraphBatched}) adds a batch dimension to all edge arrays,
achieving $30$--$40\times$ speedup over sequential decoding.


% ══════════════════════════════════════════════════════════════════════════════
\section{Neural Decoders}
\label{sec:neural}

\subsection{Neural Computational Graph (NCG) Architecture}

The neural decoder (\texttt{neural/neural\_comp\_graph.py}) replaces the
analytical tree operations with learned MLP-based operations while
preserving the $O(N \log N)$ computational graph structure.

\subsubsection{Architecture Components}

\begin{itemize}
  \item \textbf{Channel Embedding:} An \texttt{nn.Embedding} layer maps
        discrete channel outputs $z \in \{0, 1, 2\}$ (for BEMAC) to
        $d$-dimensional vectors. Root embeddings are bit-reversed before
        tree processing.

  \item \textbf{Neural calcLeft/calcRight:} Each is a shared MLP
        $\mathbb{R}^{3d} \to \mathbb{R}^d$ with ELU activations and
        configurable depth. The input concatenates the parent's first half,
        parent's second half, and the sibling edge embedding.

  \item \textbf{Decision Head (\texttt{emb2logits}):} An MLP
        $\mathbb{R}^d \to \mathbb{R}^4$ producing joint logits over
        $(u, v) \in \{0,1\}^2$. Shared between leaf decisions and
        CalcParent.

  \item \textbf{CalcParent:} A GRU-style gated residual module
        (\texttt{NeuralCalcParent}) that takes left and right child
        embeddings $\mathbb{R}^d \times \mathbb{R}^d \to \mathbb{R}^d$,
        using a gate network (sigmoid) and candidate network (ELU MLP)
        with an averaging residual connection. This fully neural approach
        has $O(md)$ complexity per operation, independent of channel structure.

  \item \textbf{Learnable ``no info'' embedding:} A learnable parameter
        vector initializes leaf edges before they receive information.
\end{itemize}

\subsubsection{Configuration}

The default configuration is $d = 16$, hidden $= 64$, 2 hidden layers per
MLP, temperature $\tau = 1.0$. The total parameter count is modest:
typically $\sim$15K--30K parameters.

\subsection{Memory Channel Extension (\texttt{ncg\_memory.py})}

The \texttt{NCGMemoryDecoder} extends the NCG to channels with memory by
replacing the discrete embedding lookup with a \textbf{Sequence Encoder}:

\begin{itemize}
  \item A bidirectional GRU processes the continuous channel output sequence
        $z_1, \ldots, z_N$, producing context-aware representations.
  \item A projection layer maps GRU outputs to $d$-dimensional root embeddings.
  \item The tree operations remain \emph{exactly the same} as the BEMAC model.
\end{itemize}

Two memory channel models are implemented in \texttt{channels\_memory.py}:
\begin{itemize}
  \item \textbf{ISI MAC:} $z_t = s_x(t) + s_y(t) + \alpha [s_x(t{-}1) + s_y(t{-}1)] + w_t$ with BPSK modulation, state space $S = 4$.
  \item \textbf{Gilbert--Elliott MAC:} 2-state Markov channel with different noise levels per state.
\end{itemize}

The key advantage is that the neural decoder's complexity is $O(md \cdot N \log N)$, independent of the channel state space size $S$, whereas analytical SC has complexity $O(S^3 \cdot N \log N)$.

\subsection{Training Procedure (\texttt{neural/train.py})}

Training uses teacher forcing with 4-class cross-entropy loss:
\begin{enumerate}
  \item Generate random $(u, v)$ pairs, encode, and transmit through BEMAC
  \item Run the NCG decoder forward pass with teacher-forced true values
  \item At each information bit position, the model outputs 4-class logits
        for $(u_t, v_t) \in \{(0,0), (0,1), (1,0), (1,1)\}$
  \item Compute cross-entropy loss and backpropagate
\end{enumerate}

The training script supports three phases:
\begin{enumerate}
  \item \textbf{Overfit test:} Validate architecture on 100 fixed samples
  \item \textbf{Full training:} Train on Class B path with Adam optimizer and cosine annealing
  \item \textbf{Evaluation:} Compare NN BLER vs.\ SC BLER
\end{enumerate}

\subsection{Saved Models}

Pre-trained models are stored in \texttt{saved\_models/}:
\begin{table}[h]
\centering
\caption{Pre-trained NCG model inventory}
\begin{tabular}{lll}
\toprule
\textbf{File} & \textbf{Block Length} & \textbf{Notes} \\
\midrule
\texttt{ncg\_N8\_d16.pt}  & $N = 8$ & Base model \\
\texttt{ncg\_N16\_d16.pt} & $N = 16$ & Base model \\
\texttt{ncg\_N32\_curriculum.pt} & $N = 32$ & Curriculum learning \\
\texttt{ncg\_N64\_best.pt} & $N = 64$ & Best validation \\
\texttt{ncg\_N128\_curriculum.pt} & $N = 128$ & Curriculum learning \\
\texttt{ncg\_N256\_bler\_sweep.pt} & $N = 256$ & BLER sweep \\
\texttt{ncg\_N512\_bler\_sweep.pt} & $N = 512$ & BLER sweep \\
\texttt{ncg\_N1024\_bler\_sweep.pt} & $N = 1024$ & BLER sweep \\
\bottomrule
\end{tabular}
\end{table}


% ══════════════════════════════════════════════════════════════════════════════
\section{Evaluation Framework}
\label{sec:eval}

\subsection{MACEval Class (\texttt{eval.py})}

The \texttt{MACEval} class provides a unified Monte Carlo evaluation pipeline:

\begin{lstlisting}[caption={Basic MACEval usage}]
from polar import BEMAC, MACEval, design_bemac, make_path

channel = BEMAC()
N, n = 1024, 10
b = make_path(N, path_i=N)
Au, Av, frozen_u, frozen_v, z_u, z_v = design_bemac(n, ku=307, kv=614)

evaluator = MACEval(channel, log_domain=True, backend='auto')
ber_u, ber_v, bler = evaluator.run(
    N, b, Au, Av, frozen_u, frozen_v, n_codewords=50000)
\end{lstlisting}

\paragraph{Backend selection:}
The \texttt{backend} parameter controls decoder dispatch:
\begin{itemize}
  \item \texttt{'auto'}: Selects the fastest decoder for the given path (efficient for extreme paths, interleaved for others).
  \item \texttt{'reference'}: Forces the $O(N^2)$ reference decoder.
  \item \texttt{'efficient'}: Forces the efficient $O(N \log N)$ decoder (extreme paths only).
  \item \texttt{'interleaved'}: Forces the computational graph decoder.
\end{itemize}

\paragraph{Adaptive MC evaluation:}
The \texttt{adaptive\_mc\_eval} function automatically scales the number
of Monte Carlo trials to fit within a given time budget by benchmarking
single-decode latency first.

\subsection{Simulation Scripts}

The \texttt{scripts/} directory contains simulation drivers (see Section~\ref{sec:scripts}).


% ══════════════════════════════════════════════════════════════════════════════
\section{Results}
\label{sec:results}

This section presents simulation results extracted from the project's JSON
result files. All BLER figures are obtained via Monte Carlo simulation with
at least 50{,}000 codewords per data point (unless otherwise noted).

\subsection{BEMAC Results: BLER vs.\ Block Length}

\subsubsection{Class C (Path $0^N 1^N$)}

Table~\ref{tab:bemac_classC} shows the BLER performance for BEMAC Class C
(path $0^N 1^N$, SC decoding, rate ratio $\rho = 0.6$). The BLER decreases
steadily with $N$, confirming polarization. The V-channel BER is
consistently lower than the U-channel BER because for Class C the
V-conditional channel is noiseless ($Z_0^{(V)} = 0$) and all V errors come
from error propagation in the U phase.

\begin{table}[h]
\centering
\caption{BEMAC Class C, SC ($L = 1$), $\rho = 0.6$, 50{,}000 codewords, analytical design}
\label{tab:bemac_classC}
\begin{tabular}{rrrrrrr}
\toprule
$N$ & $k_u$ & $k_v$ & $R_u$ & $R_v$ & BLER & BER$_u$ \\
\midrule
   8 &   2 &    5 & 0.250 & 0.625 & 0.103    & $8.3 \times 10^{-2}$ \\
  16 &   5 &   10 & 0.313 & 0.625 & 0.086    & $4.5 \times 10^{-2}$ \\
  32 &  10 &   19 & 0.313 & 0.594 & 0.087    & $3.3 \times 10^{-2}$ \\
  64 &  19 &   38 & 0.297 & 0.594 & 0.057    & $2.0 \times 10^{-2}$ \\
 128 &  38 &   77 & 0.297 & 0.602 & 0.024    & $7.0 \times 10^{-3}$ \\
 256 &  77 &  154 & 0.301 & 0.602 & 0.015    & $3.7 \times 10^{-3}$ \\
 512 & 154 &  307 & 0.301 & 0.600 & 0.003    & $8.1 \times 10^{-4}$ \\
1024 & 307 &  614 & 0.300 & 0.600 & $2.6 \times 10^{-4}$ & $5.2 \times 10^{-5}$ \\
\bottomrule
\end{tabular}
\end{table}

The rates converge to $(R_u, R_v) \approx (0.30, 0.60)$ as $N$ grows.
This is consistent with the target sum rate $R_u + R_v = 0.9$ at ratio
$\rho = R_u / (R_u + R_v) = 1/3$. The three-order-of-magnitude drop in
BLER from $N = 8$ to $N = 1024$ demonstrates clean polarization behavior.

Table~\ref{tab:bemac_classA} shows BEMAC Class A results at $\rho = 0.6$.
The BLER \emph{increases} with $N$, reaching 1.0 at $N \geq 512$. This
confirms that the symmetric rate direction is capacity-suboptimal for the
asymmetric BEMAC capacity region.

\begin{table}[h]
\centering
\caption{BEMAC Class A, SC ($L = 1$), $\rho = 0.6$, design: analytical}
\label{tab:bemac_classA}
\begin{tabular}{rrrrrr}
\toprule
$N$ & $k_u$ & $k_v$ & $R_u = R_v$ & BLER \\
\midrule
   8 &   4 &   4 & 0.500 & 0.162 \\
  16 &   7 &   7 & 0.438 & 0.273 \\
  32 &  14 &  14 & 0.438 & 0.462 \\
  64 &  29 &  29 & 0.453 & 0.764 \\
 128 &  58 &  58 & 0.453 & 0.938 \\
 256 & 115 & 115 & 0.449 & 0.997 \\
 512 & 230 & 230 & 0.449 & 1.000 \\
1024 & 461 & 461 & 0.450 & 1.000 \\
\bottomrule
\end{tabular}
\end{table}

\subsection{BEMAC Class B Results}

Table~\ref{tab:bemac_classB} shows BEMAC Class B (path $0^{N/2} 1^N 0^{N/2}$).
The BLER also increases with $N$ at $\rho = 0.6$, though less dramatically
than Class A. At small $N$ (e.g., $N = 8$), Class B achieves $\text{BLER} = 0$.

\begin{table}[h]
\centering
\caption{BEMAC Class B, SC ($L = 1$), $\rho = 0.6$, design: analytical}
\label{tab:bemac_classB}
\begin{tabular}{rrrrrr}
\toprule
$N$ & $k_u$ & $k_v$ & BLER & BER$_u$ \\
\midrule
   8 &   3 &   4 & 0.000  & 0.000 \\
  16 &   6 &   8 & 0.055  & $2.8 \times 10^{-2}$ \\
  32 &  12 &  17 & 0.542  & 0.238 \\
  64 &  24 &  34 & 0.762  & 0.323 \\
 128 &  48 &  67 & 0.884  & 0.373 \\
 256 &  96 & 134 & 0.986  & 0.411 \\
 512 & 192 & 269 & 1.000  & 0.412 \\
1024 & 384 & 538 & 1.000  & 0.410 \\
\bottomrule
\end{tabular}
\end{table}

\subsection{Gaussian MAC Results}

The Gaussian MAC campaign covers Classes A, B, C at rate ratios
$\rho \in \{0.3, 0.5, 0.7, 0.9\}$ and block lengths $N = 8, \ldots, 1024$,
using Monte Carlo genie-aided code design. The campaign data in
\texttt{results/gmac\_final\_data.json} contains over 200 data points.

\subsubsection{SC Decoder Results ($L = 1$)}

Table~\ref{tab:gmac_sc} shows selected SC BLER results extracted from the
campaign data.

\begin{table}[h]
\centering
\caption{Gaussian MAC SC ($L = 1$) BLER: selected results from the campaign}
\label{tab:gmac_sc}
\begin{tabular}{llrrrrrr}
\toprule
\textbf{Class} & $\rho$ & $N = 8$ & $N = 32$ & $N = 64$ & $N = 128$ & $N = 256$ & $N = 1024$ \\
\midrule
C & 0.3 & 0.123 & 0.043 & 0.012 & 0.001 & 0.008 & 0.012 \\
C & 0.5 & 0.109 & 0.076 & 0.027 & 0.007 & 0.002 & 0.002 \\
C & 0.7 & 0.226 & 0.139 & 0.120 & 0.113 & 0.081 & 0.024 \\
C & 0.9 & 0.198 & 0.326 & 0.492 & 0.573 & 0.652 & 0.774 \\
\midrule
A & 0.3 & 0.124 & 0.007 & 0.018 & 0.004 & 0.008 & 0.013 \\
A & 0.5 & 0.046 & 0.008 & 0.003 & 0.000 & 0.001 & 0.000 \\
A & 0.7 & 0.081 & 0.071 & 0.171 & 0.240 & 0.238 & 0.465 \\
A & 0.9 & 0.684 & 0.955 & 0.999 & 1.000 & 1.000 & 1.000 \\
\midrule
B & 0.3 & 0.067 & 0.008 & 0.002 & 0.001 & 0.005 & 0.033 \\
B & 0.5 & 0.038 & 0.007 & 0.007 & 0.000 & 0.000 & 0.001 \\
B & 0.7 & 0.071 & 0.042 & 0.031 & 0.013 & 0.004 & 0.002 \\
B & 0.9 & 0.249 & 0.505 & 0.650 & 0.726 & 0.892 & 0.998 \\
\bottomrule
\end{tabular}
\end{table}

\subsubsection{SCL Decoder Results ($L = 32$)}

Table~\ref{tab:gmac_scl} shows SCL results at $L = 32$ for selected
configurations. The list decoder provides significant gains at moderate
rates.

\begin{table}[h]
\centering
\caption{Gaussian MAC SCL ($L = 32$) BLER: selected results}
\label{tab:gmac_scl}
\begin{tabular}{llrrrrrr}
\toprule
\textbf{Class} & $\rho$ & $N = 8$ & $N = 32$ & $N = 64$ & $N = 128$ & $N = 256$ & $N = 1024$ \\
\midrule
C & 0.5 & 0.116 & 0.038 & 0.009 & 0.000 & 0.000 & 0.000 \\
C & 0.7 & 0.222 & 0.148 & 0.087 & 0.031 & 0.016 & 0.000 \\
A & 0.5 & 0.033 & 0.004 & 0.004 & 0.000 & 0.000 & 0.000 \\
A & 0.7 & 0.044 & 0.056 & 0.081 & 0.141 & 0.108 & 0.180 \\
B & 0.5 & 0.032 & 0.003 & 0.004 & 0.000 & 0.000 & 0.000 \\
B & 0.7 & 0.036 & 0.025 & 0.009 & 0.004 & 0.002 & 0.000 \\
\bottomrule
\end{tabular}
\end{table}

\subsubsection{Key Observations}

\begin{enumerate}
  \item \textbf{Rate ratio sensitivity:} At moderate rate ratios
        ($\rho \leq 0.5$), BLER decreases with $N$ for all classes, confirming
        that the codes approach capacity. At $\rho = 0.5$ for Class A,
        $\text{BLER} = 0.000$ is achieved already at $N = 128$.

  \item \textbf{High rate ratio floor:} At $\rho = 0.9$, BLER increases
        with $N$ for all classes, reaching 1.0 at large $N$. This indicates
        operation beyond the achievable rate region at the given SNR.

  \item \textbf{SCL improvement:} The SCL decoder ($L = 32$) substantially
        improves performance at $\rho = 0.7$. For example, Class B at
        $N = 1024$ goes from $\text{BLER} = 0.002$ (SC) to $0.000$ (SCL).
        At Class C with $\rho = 0.7$, SCL achieves $\text{BLER} = 0$ at
        $N = 1024$ compared to SC's $0.024$.

  \item \textbf{Class comparison:} Class B with $\rho = 0.7$ scales better
        than Class A with the same $\rho$. At $N = 1024$, Class B achieves
        $\text{BLER} = 0.002$ while Class A has $0.465$ under SC.

  \item \textbf{Non-monotone BLER:} Some configurations show non-monotone
        BLER vs.\ $N$ behavior (e.g., Class A at $\rho = 0.3$ with
        $\text{BLER}$ dipping to $0.004$ at $N = 128$ then rising to $0.013$
        at $N = 1024$). This is attributed to the discrete rate quantization:
        the actual rates $(k_u/N, k_v/N)$ cannot exactly match the target
        at every $N$.
\end{enumerate}

\subsection{ABNMAC Results}

ABNMAC simulations span Classes A, B, C with MC-based design at block lengths
$N = 8$ to $N = 1024$. The ABNMAC channel has asymmetric capacities
($I(Z;X) \approx 0.400$, $I(Z;Y|X) \approx 0.800$), making it more
challenging than BEMAC.

At $N = 8$, Class C achieves $\text{BLER} \approx 0.119$ with $k_u = 2$,
$k_v = 3$, using the $O(N \log N)$ efficient decoder with MC design.
The ABNMAC requires the Monte Carlo design procedure since the
Bhattacharyya parameters depend on the full noise distribution.

\subsubsection{ABNMAC Design Considerations}

A key difference from BEMAC is that neither user has a noiseless conditional
channel. The V-conditional Bhattacharyya starting value is
$Z_0^{(V)} \approx 0.337$ (compared to 0 for BEMAC), meaning the V-channel
also experiences polarization and not all V positions can carry information.
This reduces the achievable rates compared to the capacity bounds,
particularly at short block lengths.

\subsection{SC vs.\ SCL Comparison on BEMAC}

To quantify the benefit of list decoding, we compare SC ($L = 1$) and SCL
($L = 32$) on BEMAC for various rate ratios. The BEMAC is an ideal
testbed because the V-conditional channel is noiseless for Class C,
isolating the impact of list decoding on the U-phase.

At aggressive rate ratios ($\rho = 0.6$--$0.8$), SCL provides notable
improvements over SC, particularly at intermediate block lengths
($N = 32$--$256$) where the finite-length penalty is most pronounced.
At $\rho = 0.5$ (the most conservative setting), both SC and SCL achieve
very low BLER, making the list decoder's advantage less visible.

\subsection{Results Directory Structure}

The \texttt{results/} directory is organized as follows:
\begin{itemize}
  \item \textbf{JSON files:} Primary simulation data, named by convention
        \texttt{sim\_\{channel\}\_\{class\}\_\{design\}\_L\{list\_size\}\_\{params\}.json}.
        Each file contains metadata (channel, class, decoder type, design method)
        and a list of result records with $N$, $k_u$, $k_v$, BLER, BER$_u$,
        BER$_v$, block error count, codeword count, and wall-clock time.

  \item \textbf{Campaign files:} \texttt{campaign\_10h.json} and
        \texttt{campaign\_28h.json} contain multi-day Gaussian MAC simulation
        campaigns with hundreds of data points.

  \item \textbf{Plot files:} PNG and PDF figures generated by the plotting
        scripts, organized into \texttt{final\_plots/} (publication-quality)
        and \texttt{clean\_plots/} (intermediate) subdirectories.

  \item \textbf{NCG results:} \texttt{ncg\_vs\_sc\_classB.json} and
        \texttt{nn\_vs\_sc\_bemac\_*.json} contain neural decoder evaluation
        results.
\end{itemize}

\subsection{Generated Plots}

The plotting scripts generate the following figure types:
\begin{enumerate}
  \item \textbf{BLER vs.\ $N$ by class:} Compares Classes A, B, C at fixed
        SNR and rate ratio (e.g., \texttt{01\_classes\_snr3dB.png}).
  \item \textbf{BLER vs.\ $N$ by rate ratio:} Compares $\rho = 0.3, 0.5, 0.7, 0.9$
        for a fixed class (e.g., \texttt{02\_rho\_snr6dB.png}).
  \item \textbf{SC vs.\ SCL comparison:} Overlay SC and SCL curves
        (e.g., \texttt{03\_scl\_snr3dB.png}).
  \item \textbf{Rate vs.\ BLER:} BLER as a function of sum rate at fixed $N$
        (e.g., \texttt{04\_rate\_snr6dB.png}).
  \item \textbf{SNR sweep:} BLER vs.\ SNR for fixed $N$ and class
        (e.g., \texttt{05\_snr\_compare.png}).
  \item \textbf{Heatmap:} BLER as a function of $(N, \rho)$ for a given
        class and SNR (e.g., \texttt{05\_heatmap\_clean.png}).
\end{enumerate}

\subsection{NCG Neural Decoder Results}

Table~\ref{tab:ncg_vs_sc} compares the NCG neural decoder against analytical
SC on BEMAC Class B.

\begin{table}[h]
\centering
\caption{NCG vs.\ SC BLER on BEMAC Class B}
\label{tab:ncg_vs_sc}
\begin{tabular}{rrrrr}
\toprule
$N$ & $k_u$ & $k_v$ & NN BLER & SC BLER \\
\midrule
   8 &  4 &  6 & 0.0586 & 0.0582 \\
  16 &  8 & 11 & 0.0106 & 0.0118 \\
  32 & 16 & 22 & 0.0082 & 0.0096 \\
  64 & 32 & 45 & 0.0044 & 0.0056 \\
 128 & 64 & 90 & 0.0018 & 0.0034 \\
\bottomrule
\end{tabular}
\end{table}

The neural decoder matches SC at $N = 8$ and progressively outperforms it
as $N$ increases. At $N = 128$, the NCG achieves a $47\%$ reduction in
BLER ($0.0018$ vs.\ $0.0034$). This is achieved with curriculum learning
and a modest model size ($d = 16$, hidden $= 64$).


% ══════════════════════════════════════════════════════════════════════════════
\section{Implementation Notes}
\label{sec:implementation}

This section documents key implementation details and design decisions that
are important for understanding and extending the codebase.

\subsection{Indexing Conventions}

The project consistently uses \textbf{1-based indexing} for polar code
positions:
\begin{itemize}
  \item Information sets $A_U$, $A_V$ contain 1-indexed positions
  \item Frozen dictionaries map 1-indexed positions to frozen values
  \item The decoder's internal coordinate probabilities use 1-indexed $(i, j)$
  \item The computational graph uses 1-indexed edges ($1, \ldots, 2N{-}1$)
        and vertices ($1, \ldots, N{-}1$)
\end{itemize}

Internal arrays (e.g., Bhattacharyya parameters, error rates) use 0-indexed
NumPy arrays. The conversion between conventions is consistently handled at
the interface boundaries.

\subsection{Log-Domain Arithmetic}

All probability computations are performed in the log-domain to avoid
numerical underflow. The key operations are:
\begin{itemize}
  \item \texttt{np.logaddexp(a, b)}: Computes $\log(e^a + e^b)$ stably
  \item Products become sums: $\log(p \cdot q) = \log p + \log q$
  \item Normalization: subtract $\log \sum_i e^{x_i}$ from all entries
  \item Zero probability: represented as $-\infty$ (\texttt{-np.inf})
\end{itemize}

The circular convolution involves a 4-term \texttt{logaddexp} for each of
the 4 output entries, totaling 16 additions and 12 \texttt{logaddexp}
operations per tensor pair. This is the computational bottleneck of the
decoder.

\subsection{Numerical Stability}

Special care is taken for edge cases:
\begin{itemize}
  \item \textbf{NaN handling in LLR f-node:} When both $\Lambda_a$ and
        $\Lambda_b$ are $\pm\infty$, the \texttt{logaddexp} formula can
        produce NaN. The implementation falls back to the min-sum
        approximation: $f(\Lambda_a, \Lambda_b) \approx \text{sign}(\Lambda_a)
        \cdot \text{sign}(\Lambda_b) \cdot \min(|\Lambda_a|, |\Lambda_b|)$.
  \item \textbf{Normalization of all-$-\infty$ tensors:} The
        \texttt{norm\_prod} function checks whether the total is finite
        before subtracting; if not, the tensor is left unnormalized.
  \item \textbf{BEMAC deterministic transitions:} For BEMAC, the leaf
        log-probabilities are exactly $0$ or $-\infty$, making the
        arithmetic exact (no floating-point approximation).
\end{itemize}

\subsection{Performance Optimization}

The project employs several levels of optimization:

\begin{enumerate}
  \item \textbf{Vectorized batch encoding:} The \texttt{polar\_encode\_batch}
        function processes all codewords simultaneously using NumPy array
        operations, with the butterfly XOR stage using reshape + slice + XOR.

  \item \textbf{Batched SC decoding:} The \texttt{\_CompGraphBatched} class
        adds a batch dimension to all edge arrays. All batch elements share
        the same tree traversal pattern (since they use the same path $b$),
        so the control flow is identical---only the data differs. This
        achieves $30$--$40\times$ speedup over sequential decoding.

  \item \textbf{Numba JIT compilation:} The parallel decoder
        (\texttt{decoder\_parallel.py}) compiles all inner loops to native
        code via Numba's \texttt{@njit} decorator. The flat array storage
        avoids Python object overhead. The \texttt{prange} directive enables
        automatic parallelization across CPU cores for batch decoding.

  \item \textbf{Lazy evaluation:} The \texttt{\_SCNode} tree only computes
        LLRs when requested (\texttt{get\_llr(k)}), and the computational
        graph only evaluates edges when navigated to (\texttt{step\_to}).
\end{enumerate}

\subsection{Numba Closure Bug}

An important lesson from the parallel decoder development: Numba closures
that capture mutable arrays (e.g., the edge data array) can produce
\emph{silently incorrect results} because Numba may not properly track
mutations to captured arrays. The solution is to make all \texttt{@njit}
functions \emph{standalone}, taking the mutable array as an explicit
argument rather than capturing it from a closure scope.

\subsection{Design Comparison: Analytical vs.\ Monte Carlo}

The project provides a \texttt{compare\_with\_analytical} utility that runs
both the analytical Bhattacharyya design and the MC design on the same
channel and reports whether the resulting information sets agree. For BEMAC
with the extreme path $0^N 1^N$, the two methods produce identical designs
at all tested block lengths, validating both implementations.

For ABNMAC and Gaussian MAC, the MC design is the primary method since the
analytical approach requires closed-form starting parameters that may not
capture the full complexity of intermediate paths.


% ══════════════════════════════════════════════════════════════════════════════
\section{Scripts Reference}
\label{sec:scripts}

\begin{longtable}{p{5.5cm}p{9.5cm}}
\toprule
\textbf{Script} & \textbf{Description} \\
\midrule
\endfirsthead
\toprule
\textbf{Script} & \textbf{Description} \\
\midrule
\endhead
\texttt{scripts/simulate.py}              & General-purpose SC/SCL simulator for BEMAC \\
\texttt{scripts/simulate\_abnmac.py}      & ABNMAC-specific SC simulator \\
\texttt{scripts/simulate\_gmac.py}        & Gaussian MAC simulator with SNR parameter \\
\texttt{scripts/bler\_vs\_N.py}           & Sweep BLER vs.\ $N$ for a given channel and class \\
\texttt{scripts/bler\_vs\_N\_bemac.py}    & BEMAC-specific BLER vs.\ $N$ sweep with analytical design \\
\texttt{scripts/run\_bler\_vs\_N.py}      & Batch runner for BLER vs.\ $N$ across classes \\
\texttt{scripts/run\_design.py}           & Run MC design for BEMAC and save to \texttt{designs/} \\
\texttt{scripts/run\_design\_abnmac.py}   & Run MC design for ABNMAC \\
\texttt{scripts/run\_design\_gmac.py}     & Run MC design for Gaussian MAC \\
\texttt{scripts/campaign\_10h.py}         & 10-hour Gaussian MAC simulation campaign \\
\texttt{scripts/campaign\_28h.py}         & 28-hour extended campaign \\
\texttt{scripts/overnight\_gmac.py}       & Overnight Gaussian MAC campaign \\
\texttt{scripts/benchmark\_parallel.py}   & Benchmark Numba parallel decoder speedup \\
\texttt{scripts/demo\_one\_codeword.py}   & Single codeword encode/decode demonstration \\
\texttt{scripts/plot\_results.py}         & General BEMAC result plotter \\
\texttt{scripts/plot\_abnmac.py}          & ABNMAC result plotter \\
\texttt{scripts/plot\_bler\_vs\_N\_all\_classes.py} & Multi-class BLER vs.\ $N$ comparison plots \\
\texttt{scripts/plot\_gmac\_final.py}     & Gaussian MAC final publication-quality plots \\
\texttt{scripts/generate\_decoder\_pdf.py}& Generate decoder architecture diagrams \\
\bottomrule
\end{longtable}


% ══════════════════════════════════════════════════════════════════════════════
\section{File Structure and API Reference}
\label{sec:filestructure}

\subsection{Project File Tree}

\begin{verbatim}
to_git_v2/
  polar/                        # Core library
    __init__.py                 # Public API exports
    channels.py                 # BEMAC, ABNMAC, GaussianMAC
    encoder.py                  # Polar encoder (BR + butterfly)
    decoder.py                  # Unified SC decoder (auto-dispatch)
    decoder_scl.py              # SC List decoder
    decoder_interleaved.py      # O(N log N) comp. graph decoder
    decoder_parallel.py         # Numba JIT parallel decoder
    efficient_decoder.py        # O(N log N) LLR decoder (extreme)
    _decoder_numba.py           # Low-level Numba helpers
    design.py                   # Analytical Bhattacharyya design
    design_mc.py                # Monte Carlo genie-aided design
    eval.py                     # MACEval evaluation pipeline

  neural/                       # Neural decoder
    __init__.py
    neural_comp_graph.py        # Baseline NCG decoder
    ncg_pure_neural.py          # Pure Neural CalcParent decoder (38K params)
    ncg_memory.py               # NCG with GRU sequence encoder
    channels_memory.py          # ISI MAC, Gilbert-Elliott MAC
    train_pure_neural.py        # Training script

  designs/                      # Pre-computed design files (.npz)
  saved_models/                 # Pre-trained NCG models (.pt)
  results/                      # Simulation results (.json, .png)
    final_plots/                # Publication-quality figures
    clean_plots/                # Cleaned figures
  scripts/                      # Simulation and plotting scripts
  docs/                         # Documentation
  design.py                     # Top-level design entry point
\end{verbatim}

\subsection{Public API (\texttt{polar/\_\_init\_\_.py})}

The public API exports the following:

\begin{table}[h]
\centering
\caption{Public API}
\begin{tabular}{lll}
\toprule
\textbf{Module} & \textbf{Symbols} & \textbf{Purpose} \\
\midrule
\texttt{encoder} & \texttt{polar\_encode}, \texttt{build\_message} & Encoding \\
\texttt{channels} & \texttt{BEMAC}, \texttt{ABNMAC}, \texttt{GaussianMAC} & Channels \\
\texttt{design} & \texttt{design\_bemac}, \texttt{make\_path} & Design \\
\texttt{design\_mc} & \texttt{design\_bemac\_mc} & MC design \\
\texttt{decoder} & \texttt{decode\_single}, \texttt{decode\_batch} & SC decoding \\
\texttt{decoder\_scl} & \texttt{decode\_single\_list}, \texttt{decode\_batch\_list} & SCL decoding \\
\bottomrule
\end{tabular}
\end{table}

Additional internal modules:
\begin{itemize}
  \item \texttt{design\_mc.design\_from\_file()}: Load MC design from \texttt{.npz} and build info/frozen sets.
  \item \texttt{eval.MACEval}: Full Monte Carlo evaluation pipeline with backend selection.
  \item \texttt{eval.adaptive\_mc\_eval()}: Time-budgeted evaluation.
\end{itemize}


% ══════════════════════════════════════════════════════════════════════════════
\section{References}
\label{sec:references}

\begin{thebibliography}{9}

\bibitem{arikan2009}
E.~Arikan, ``Channel polarization: A method for constructing
capacity-achieving codes for symmetric binary-input memoryless channels,''
\emph{IEEE Trans.\ Inform.\ Theory}, vol.~55, no.~7, pp.~3051--3073, July 2009.

\bibitem{onay2013}
S.~\"{O}nay, ``Successive cancellation decoding of polar codes for the
two-user binary-input MAC,''
in \emph{Proc.\ IEEE Int.\ Symp.\ Inform.\ Theory (ISIT)}, Istanbul, Turkey,
July 2013, pp.~1368--1372.

\bibitem{ren2025}
Y.~Ren, R.~Nasser, and U.~Bhatt, ``Polar codes for the multiple access
channel via monotone chain decoding,''
2025.

\bibitem{talvardy2013}
I.~Tal and A.~Vardy, ``How to construct polar codes,''
\emph{IEEE Trans.\ Inform.\ Theory}, vol.~59, no.~10, pp.~6562--6582,
Oct.\ 2013.

\end{thebibliography}

\end{document}
